\section{Introduction and task specification}
\label{sec:intro}
Federal HR rules are spread across Title 5 of the U.S.\ Code, Title 5 of the Code of Federal Regulations (eCFR)
and OPM fact sheets. They are written in nested legal paragraphs such as \texttt{(a)(2)(i)}, and the answers often
need a table lookup or arithmetic, for example GS pay by locality or leave accrual by years of service. A useful assistant
must find the right paragraph, compute rather than guess, cite its source, and know when to stop.

\textbf{Users.} Federal employees and HR specialists asking about \emph{rules} (not about their personal records).

\textbf{What it does.} (i) Answers policy questions (annual, sick, court, military, FMLA, bereavement, weather and
safety leave; premium pay; hours of duty; incentives; within-grade increases) from the corpus with clause-level
citations. (ii) Computes leave accrual and carryover caps, looks up 2026 GS salary, hourly and overtime rates, and
lists federal holidays, using deterministic tools. (iii) Drafts an HR case for anything that needs a human decision.

\textbf{What it refuses.} Another named person's pay, leave or records (privacy); topics outside US federal HR
(private-company benefits, investment advice); requests to reveal its instructions or to change role (prompt
injection). It never submits an HR case itself and never decides an individual's eligibility.

\textbf{What a correct answer is.} Each eval item specifies the \emph{expected behaviour} (\texttt{answer},
\texttt{refuse}, \texttt{escalate}, \texttt{not\_covered}) and \emph{required facts} as groups of accepted phrasings
(e.g.\ the groups \{``240''\} and \{``hour''\} for the carryover cap). An answer is correct when the behaviour matches and every fact group is matched.
We also record whether a verified citation is present, whether a gold clause was retrieved, and which tools ran.
When the corpus lacks the answer, the system must reply ``The provided policy documents do not cover this.''
